\subsection{半范数的定义}

定义 1. --- 设 E 是 K 上的一个左向量空间。E 到 \(R_{+} = \left[0, +\infty\right]\) 中的映射 p 称为 E 上的半范数，如果它满足下列公理：

\((\mathrm{SN}_{\mathrm{I}})\) 若 \(x \in \mathrm{E}\) 且 \(\lambda \in \mathrm{K}\)，则 \(p(\lambda x) = |\lambda| p(x)\)。

\((\mathrm{SN}_{\mathrm{II}})\) 若 \(x \in \mathrm{E}\) 且 \(y \in \mathrm{E}\)，则 \(p(x + y) \leqslant p(x) + p(y)\)。

由于 \(p(x) \leqslant p(y) + p(x - y)\) 及 \(p(y) \leqslant p(x) + p(y - x)\)，由 \(p(y - x) = p(x - y)\)，我们推出

\[
\left| p (x) - p (y) \right| \leqslant p (x - y).\tag{1}
\]

例. --- 1) E 上的范数是这样一个半范数 p：关系 \(p(x) = 0\) 蕴含 x = 0（I, 第 3 页）。

\begin{enumerate}
\def\labelenumi{\arabic{enumi})}
\setcounter{enumi}{1}
\item
  对每个线性型 \(f\)（在 \(\mathbf{E}\) 上），函数 \(x \mapsto |f(x)|\) 是 \(\mathbf{E}\) 上的半范数。
\item
  若 \(p_i (1 \leqslant i \leqslant n)\) 是 \(E\) 上有限个半范数，则显然 \(p'(x) = \sup_{1 \leqslant i \leqslant n} p_i(x)\)。
\end{enumerate}

与 \(p''(x) = \sum_{i=1}^{n} \alpha_i p_i(x)\)（其中诸 \(\alpha_i\) \(\geqslant 0\)）都是 E 上的半范数。

\(p\) 是 \(\mathbf{E}\) 到 \(\mathbf{R}_{+}\) 中的映射，若它满足 \((\mathrm{SN}_{\mathrm{I}})\) 及下列公理，则称为超半范数：

\((\mathrm{SN}_{\mathrm{II}}^{\prime})\) 若 \(x \in \mathrm{E}\) 且 \(y \in \mathrm{E}\)，则 \(p(x + y) \leqslant \sup(p(x), p(y))\)。

显然，超半范数是半范数。

说 K 上的绝对值是超度量的 (CA, VI, § 6.2)，就是指它是左向量空间 \(K_{s}\) 上一个不恒为零的超半范数。

命题 1. --- 设 E 是 K 上的一个左拓扑向量空间，p 是 E 上的一个半范数。下列条件等价：

\begin{enumerate}
\def\labelenumi{\alph{enumi})}
\item
  \(p\) 在 E 中连续。
\item
  \(p\) 在点 0 连续。
\item
  \(p\) 一致连续。
\item
  对每个实数 \(\alpha > 0\)，集 \(\mathbf{W}(p, \alpha)\)（诸 \(x \in \mathbf{E}\)（满足 \(p(x) < \alpha\)）所组成的）在 \(\mathbf{E}\) 中是开的。
\item
  存在实数 \(\alpha > 0\)，使得 \(\mathrm{W}(p, \alpha)\) 是 E 中 0 的一个邻域。
\item
  对每个实数 \(\alpha > 0\)，集 \(\mathrm{V}(p, \alpha)\)（诸 \(x \in \mathrm{E}\)（满足 \(p(x) \leqslant \alpha\)）所组成的）是 \(\mathrm{E}\) 中 0 的一个邻域。
\end{enumerate}

事实上，蕴含关系 \(c) \Rightarrow a) \Rightarrow b) \Rightarrow d) \Rightarrow e) \Rightarrow f) \Rightarrow c)\) 由 \((\mathrm{SN}_{\mathrm{I}})\) 及不等式 (1) 立即得出。

推论. --- 若 p 是 E 上的一个连续半范数，q 是满足 \(q \leqslant p\) 的半范数，则 q 在 E 中连续。

当 \(p\) 是 \(\mathbf{E}\) 上的超半范数时，集 \(\mathrm{W}(p, \alpha)\) 与 \(\mathrm{V}(p, \alpha)\) 都是既开又闭的。因为，我们已看到 \(\mathrm{W}(p, \alpha)\) 是开的；另一方面，若 \(z\) 是 \(\mathrm{W}(p, \alpha)\) 的一个接触点，则存在 \(y \in \mathrm{W}(p, \alpha)\) 使得 \(p(y - z) < \alpha\)，而由 \((\mathrm{SN}_{\mathrm{II}}')\) 有 \(p(z) < \alpha\)，故 \(\mathrm{W}(p, \alpha)\) 是闭的。又，\(\mathrm{V}(p, \alpha)\) 是闭的（因为 \(p\) 连续）；进而，若 \(p(x) \leqslant \alpha\) 且 \(p(y) \leqslant \alpha\)，则 \(p(x + y) \leqslant \alpha\)（由 \((\mathrm{SN}_{\mathrm{II}}')\)），这就表明 \(\mathrm{V}(p, \alpha)\) 是开的。

